\chapter{Background Research}
\sloppy
This chapter reviews the current state of research and available tools in the field of drug-target interaction prediction. The review is divided into (i) foundational concepts, (ii) current machine-learning approaches, (iii) existing prediction tools, and a (iv) gap analysis highlighting the need for the proposed project.
\section{Initial concept}
Drug discovery is the process through which potential new therapeutic entities are identified, using a combination of computational, experimental, translational, and clinical models \citep{Zhou2017}.
Drug-Target Interaction (DTI) prediction is a critical aspect of drug discovery, as it helps in identifying potential drug candidates that can interact with target proteins to exert their desired therapeutic effects \citep{Yang2024}.
A drug is a chemical compounds which cause physiological changes in the body when consumed, injected or absorbed. A target, also known as a biological target, is a structure located in an organism that is recognized or bound by other substances such as ligands or drugs and can be acted upon by a drug or other targeted molecules \citep{Overington2006}.
The pharmaceutical industry faces mounting financial pressures and long development cycles in drug discovery which creates a huge challenge when discovering new drugs, especially when they are need in short time spans, as seen and experienced with COVID-19. One way to combat these challenges is to utilize machine learning to potentially automate the initial screening process, reducing both equipment and labour costs. This approach is very valuable in the initial stages of drug discovery when the sample size is extremely large and can be accurately whittled down by an algorithm trained to predict interactions based on known patterns.


\section{Foundations of DTI Machine Learning}
In the beginning of DTI using machine learning, methods included similarity-based approaches. This is where drugs with similar structure are assumed to interact with similar proteins (Find citation for this). Feature-based machine-learning models are much more refined and rely on more scientific knowledge compared to just their structure, such as fingerprints or protein sequence features. These come in 2 different ways, Regression - where the model predicts continuous binding affinity and Binary classification - where affinity values put into thresholds, active or inactive.

\section{Deep learning models in DTI}
Deep Learning has become a strong way to determine interactions, due to its ability to learn high dimensional representations from raw data. There are many approaches for deep learning architectures, for example the use of fully connected neural networks (FCNN), \citep{Talukder2025}\citep{Zhan2025} which is a layered network that connects every neuron in one layer to every neuron in the next layer and evaluated using Accuracy, Precision, Recall, AUC-ROC, and AUC-PR.

\subsection{Sequence-Based Models}
A Sequence based model encodes raw data, Drug SMILES and Protein amino acid sequences, directly. Recurrent Neural Networks (RNN) and Convolutional Neural Network (CNN) are used to learn patterns in local chemical environments or amino-acid motifs. Models like DeepDTA \citep{Ozturk2018} displays impressive predictive performance using sequence-only architectures.

\subsection{Graph-Based models}
Graph Neural Networks (GNNs) take as input the structure of molecules as graphs of atoms, learning structural information which, to sequence based models, are difficult to learn. \citep{Yang2024} utilized this as part of their algorithm using Message Passing Neural Network (MPNN) to encode drug sequences as a part of a greater system. An MPNN can extract more detailed relational features, for improved modeling of chemical structure.

\subsection{Hybrid Networks} 
A Hybrid network is the combination of multiple neural network architectures to capture complimentary inormation to improve the learning rate of the model. We will also use Yang's \citep{Yang2024} whole hybrid model, as well as an MPNN to encode drug sequences and a Convolutional Neural Network (CNN) to encode protein sequences. To generate the other prediction it uses a High Order Graph attention Convolutional Network (HOAGCN), Which makes the prediction based on the structure of the drug molecule and target protein. All these combined make for a much more meaningful prediction producing accurate results.


\section{Binary Classification in DTI}
There's lots of research and studies that have been done using binary classification for the interaction task, this is where the model outputs whether it will (1) or it won't (0). This has become increasingly more popular as it simplifies the learning objective allowing the model to spend more computational power on classification metrics.

\subsection{Research Papers Using Binary Classification}
A strong model which uses this type of classification is DeepDTA \citep{Ozturk2018}, which applies a threshold to continuous affinity values (IC50 or Ki) and classifies the interactions as active or inactive. Although the study also explored regression, it showed that the binary classification improves model stability and reduces noise introduced by inconsistent conditions.
Another strong case is produced in Talukder \citep{Talukder2025}, who uses explicit affinity binarization combined with MACCS fingerprints and protein sequence embeddings to train a neural network that outputs a binary interaction score. This is supported by Zhan \citep{Zhan2025}, who follows the same approach, emphasizing the robustness in binary prediction when affinity values are inconsistent. 

\subsection{Affinity Thresholds}
The affinity thresholds used are very important because it determines the data the model ingests and whats its classified at. Having an off centered threshold can weight the dataset to one way even if it went in 50:50. In DeepDTA \citep{Ozturk2018} they used the thresholds for evluation againts other datasets (Davis, KIBA). They used a pKd >= 7 as binding on the Davis dataset and on the KIBA dataset they used a KIBA score >= 12.1 as the threshold. This is pretty muhc used a golden standrda among the indsurty as later works reference DeepDTA for their threshold choices.

\subsection{Negative Sampling}
Some of the recent work that uses binary classification also highlight the importance of negative sampling during training. Since many DTI datasets tend to be weighted towards positive pairs, the data becomes skewed. Several recent models including DeepPurpose \citep{DeepLearning} and other BindingDB approaches, address skewed data using negative sampling to generate synthetic pairs to avoid having an unbalanced training dataset.

\subsection{Evaluation metrics}
Commonly seen among these studies in binary classification, \citep{DeepLearning, Talukder2025, Ozturk2018}, they rely (i) Accuracy, how many observations were correctly classified (ii) Precision, How many positive predictions were actually positive (iii) Recall, how many positive cases did the model correctly predict (iv) AUC for either (i)PR, A chart that visualizes the Precision and Recall in a single graph or (ii)ROC, is a chart that visualizes the trade-off between the true positive rate (TPR) and the false positive rate (FPR).

\subsection{Binary Classification Summary}
Binary classification has become the basic method for contemporary DTI research because it can complete the tasks at machine-learning benchmarks but in a more simple way. This is why i have decided to pursue a model that works with binary classification and use classification based evaluation metrics.

\section{Dataset \& Feature Engineering}
Recent research papers consistently take as input drug SMILES and proteins as Amino Acid Sequences from BindingDB's dataset \citep{Vefghi2025}, reinforacing the fact that this is the most appropriate source data in DTI. 
A few studies differ mainly in how they encode these inputs. For example the paper by Talukder\citep{Talukder2025} represented the drugs as MACCS fingerprints generated by RDKit, these are fast, interpretable and simple but are not very expressive making them easy to compute which can benefit me with limited computational power. Protein features are similarly derived from amino acid sequences through embedding frameworks such as ProtBERT.
While datasets like Chembl are available, the frequent occurrence of BindingDB in recent literature could suggest that it's the current benchmark for DTI projects.

\section{Common limitations}
What seemed to be a common limitation was the lack of data on successful binds creating an imbalanced dataset, however some of these papers tried to combat this issue by taking data from two datasets \citep{Yang2024} or manipulating the data through affinity harmonization \& binarization \citep{Talukder2025} and using Polyloss framework \citep{Leng2022_proceedings}, which is an addition to standard loss functions that allows us to precisely tune how the model handles our severe class imbalance.

\section{Current products}

\subsection{SwissTargetPrediction} \citep{Gfeller2013} \citep{STP_web} Is a website which allows the user to predict using a ligand-based approach built to combine different similarity measures, providing more accurate predictions. The method, as described in the 2013 paper, was trained on a set of 224,412 molecules known to interact with 1700 human proteins from the ChEMBL database.

For many bioactive molecules their primary target is unknown and even when it is, most have more than one target (Off-targets) that are not well characterized. Making predictions for these off-targets is important for (i)understanding side effects and (ii)drug repositioning - finding new uses for old drugs.

The method's main strength is its hybrid approach to similarity, combining multiple variables to discover unknown targets of similar structure. It uses and combines 2D Chemical Similarity: This uses FP2 chemical fingerprints to find structurally similar molecules, which is effective for identifying analogs(structurally similar molecules) within the same chemical series. (ii) 3D Shape Similarity: This uses Electroshape, which is a 3D method that compares molecules based on their shape, atomic partial charges, and lipophilicity. This is powerful for finding active molecules with different chemical structures.

These two scores are weighted and combined using a multiple logistic regression model. The authors demonstrated that while 2D fingerprints alone perform well on biased datasets, the combined 2D/3D approach is significantly more accurate in the more realistic scenario of predicting targets for novel molecules that do not belong to a known chemical series.

The interface is plain and a bit scattered but provides an easy use, you select a species (i)Homo sapiens, (ii)Mus musculus, (iii)Rattus norvegicus, and then choose a drug molecule through given example, SMILES representation, or drawing the molecule. which is creative and more advanced way to really test new drug compounds and structures. 



\subsection{DINIES} \citep{Yamanishi2014} \citep{DINIES_web} (Drug-target Interaction Network Inference Engine based on Supervised Analysis) is a web tool that predicts potential interactions based on drug data and omics-scale protein data, allowing the user to use whatever data as input as long as it's represented accordingly. First published in 2014 it uses machine learning and predicts using the "guilt-by-association" principle, similar drugs are likely to interact with similar proteins.

The methodology is based on using Kernels to convert (i)chemical structure (ii) side-effect and (iii) Protein sequence profiles to separate matrices. The server combines the matrices and in its simplest mode will just average them together to create a single combined drug similarity matrix.

However this foundational machine learning method has become very limited compared to other products as it over relies on features and it only learns simple mapping, not complex mapping or non-linear relationships from raw data. Furthermore, the UI is not designed for simple queries as it requires the user to format and upload their own data matrices, creating a high barrier to entry for most researchers.


\section{Gap analysis}
After a detailed review of Drug-Target interaction products and papers, from foundational similarity-based methods to modern deep learning models, three gaps in the current landscape have appeared: (i) \textbf{SOTA vs Usability Gap}, gap between modern, cutting-edge research models and tools accessible to the public (ii) \textbf{Novelty problem}, similarity-based models will struggle when given novel drugs or targets which have no known analogs (iii) \textbf{The Data Imbalance gap}, older datasets were imbalanced and some using conventional loss functions to counteract but they can turn out biased, resulting in poor performance.

My proposed predictive ML pipeline will be designed to directly address these three gaps, (i) It will deliver SOTA accuracy by using modern neural network architectures to learn representations reducing the novelty issue (ii) utilizing the PolyLoss function to force the model to focus on hard-to-find positive interactions to address the data imbalance challenge (iii) diminishing the usability gap by connecting powerful backend to a simple UI, displaying results and information in an accessible way, a feature missing from most complex models.



